% Final
\subsection{Codecademy}

% Final
Druhou analyzovanou aplikací je aplikace Codecademy\footnote{Dostupné na: \href{https://www.codecademy.com/}{https://www.codecademy.com/}}. Codecademy je jednou z nejpopulárnějších platforem na výuku programování a softwarového inženýrství. Nabízí přes 800 kurzů \cite{CodecademyCourses} a má napříč platformami přes 50 milionů uživatelů \cite{CodecademyPartnerships}. Codecademy je k dispozici jako webová aplikace a dále také jako aplikace Codecademy Go, která je k dispozici na platformách Android a iOS.

% Final
Mobilní aplikace Codecademy Go obsahuje pouze jednoduché funkce pro opakování a procvičování látky. Hlavní náplň kurzů se však nachází ve webové aplikaci Codecademy. Proto se v této analýze primárně zaměřím na webovou aplikaci, nicméně zahrnu do ní i poznatky ze zmíněné mobilní aplikace.  

% Final
\subsubsection{Vzhled a uživatelská zkušenost}

% Odstavec: Obecný dojem (komplexita/přehlednost UI, paleta barev, typografie)
% Final
Codecademy na první pohled působí mnohem komplexnějším dojmem než předchozí analyzovaná aplikace Mimo. Design uživatelského rozhraní je spíše minimalistický, nicméně místy může stále působit příliš složitě. V aplikaci je jako barva pozadí využívána primárně světle růžová barva. Při učení je pak pro pozadí využívána bílá, černá nebo tmavě šedá barva. Pro tlačítka a další důležité prvky rozhraní pak aplikace využívá vysoce kontrastní barvy jako například modrou nebo žlutou. Co se typografie týče, Codecademy využívá napříč celou aplikací, s výjimkou bloků kódu, bezpatkové písmo Apercu.

% Odstavec: Rozložení prvků na obrazovce (lekce v kurzu, učení)
% Final
Rozložení prvků na obrazovce je poměrně přehledné a logické, nicméně některé prvky rozhraní mohou být pro uživatele zbytečně matoucí nebo složité. Kupříkladu hned na hlavní stránce kurzu se nachází výrazné tlačítko pro smazání pokroku celého kurzu. Přestože je tato funkce velice důležitá, nejspíše nebude tak často využívána a pozice a výraznost tohoto tlačítka může akorát mást uživatele a odpoutávat jeho pozornost.

% Final
Na hlavní stránce kurzu je obsah kurzu rozdělen do seznamu rozbalovacích modulů, v nichž se mohou nacházet různé lekce, kvízy nebo projekty. Struktura kurzu je tak uspořádána přehledně a je snadné se v ní orientovat.

% Odstavec: Navigace v aplikaci
% Final
Pro navigaci v aplikaci je využita horní navigační lišta, ve které se nachází rozbalovací menu s odkazy na všechny důležité části aplikace. Dále bývá na jednotlivých stránkách využívána levá boční navigační lišta, která slouží k navigaci mezi funkcemi, které jsou relevantní pro danou stránku.

% Odstavec: Srovnání webového a mobilního prostředí
% Final
Mobilní aplikace Codecademy Go je, co se funkcí týče, výrazně jednodušší. Umožňuje základní funkce zapisování si kurzů a učení se jednotlivých témat. V rámci kurzu jsou k dispozici pouze výrazně zjednodušené lekce, které se liší od lekcí a aktivit ve webové aplikaci.

% Odstavec: Animace a efekty a mikro interakce
% Final
V aplikaci jsou animace a efekty použity velice zřídka nebo téměř vůbec. Celá stránka tak působí velice statickým dojmem. 

% Final
\subsubsection{Způsoby učení}

% Struktura kurzu (lineární, tematické moduly, career paths, kapitoly atd.)
% Final
Co se struktury kurzů týče, Codecademy nabízí hned několik druhů kurzů. Mezi ně patří kurzy, kariérní cesty a cesty schopností. Všechny tyto druhy kurzů mají stejnou strukturu a liší se zejména svojí obsáhlostí. Dále Codecademy nabízí řadu doplňujících materiálů, kterými jsou různé nástroje pro učení, dokumentace, videa, články a další.

% Final
Jednotlivé kurzy jsou rozděleny do modulů. V nich se pak nachází různé lekce, kvízy, projekty, články nebo další aktivity.

% Jak se spustí učení a jak vypadá study session
% Final
Na hlavní stránce kurzu je uživateli zobrazeno tlačítko, které mu umožňuje spustit lekci nebo aktivitu, u které uživatel naposledy v rámci kurzu skončil. Jednotlivé lekce se pak typicky skládají ze cvičení, která zabírají přibližně 10 minut.

% Final
Cvičení se skládají z textového vysvětlení látky a dále jednotlivých úkolů. Uživatel má k dispozici textový editor, ve kterém může psát a editovat připravený kód. Dále má k dispozici okno, ve kterém je vidět výstup daného kódu.

% Final
Kromě zmíněných prvků cvičení si může uživatel také otevřít tzv. tahák se shrnutím daného tématu. Dále si může uživatel spustit přiložené YouTube video, ve kterém lektor prochází daným cvičením a vysvětluje jednotlivé kroky.

% Final
Učení je přehledné a uživatel má k dispozici dostatek materiálů pro procvičení a pochopení dané látky. Potenciální nevýhodou může být struktura samotných cvičení. Uživatel je nucen si před každým cvičením přečíst několik stránek textu, což může být pro uživatele zdlouhavé a nezáživné.

% Final
V rámci učení má uživatel k dispozici AI asistenta, který mu v chatovém okně může vysvětlit danou látku nebo poradit s řešením daného úkolu.

% Druhy cvičení
% Final
V rámci lekcí jsou uživateli zobrazována především cvičení, v nichž uživatel doplňuje a přepisuje připravený kód přímo v editoru, nebo píše kód od začátku na základě zadání.

% Final
V rámci kvízů jsou uživateli zobrazovány různé otázky, typicky formou výběru z několika odpovědí nebo výběru, zda je daný výrok pravdivý nebo nepravdivý.

% Final
Jednou z dostupných aktivit jsou takzvané projekty, které svou strukturou připomínají delší cvičení v jednotlivých lekcích. Hlavním rozdílem je, že uživatel nemá možnost nahlédnout do správného řešení a jeho řešení není nijak kontrolováno. Uživatel tak musí daný úkol vyřešit sám bez pomoci. Jednotlivé úkoly pak sám může označit za splněné.  

% Final
\subsubsection{Způsoby motivace uživatele}

% Jaké formy gamifikace aplikace používá
% Final
Na rozdíl od aplikace Mimo nepoužívá Codecademy zdaleka tolik forem gamifikace. Aplikace ovšem motivuje uživatele několika následujícími způsoby.

% Final
\paragraph{Ukazatele pokroku} Uživatel má jak na domovské stránce, tak na stránce kurzu k dispozici ukazatele pokroku, které znázorňují, kolik procent daného kurzu uživatel dokončil. 

% Final
\paragraph{Ocenění} Za splnění určitých cílů v rámci učení uživatel může získat různá ocenění, která se následně zobrazují na jeho profilu ve formě odznaků.

% Final
\paragraph{Certifikáty} Codecademy nabízí dva typy certifikátů. První typ certifikátu uživatel může získat po dokončení určitého kurzu. Druhý typ certifikátu je tzn. profesní certifikát. Ten může uživatel získat, pokud splní všechny testy v rámci dané kariérní cesty.

% Final
\paragraph{Streak} Podobně jako v aplikaci Mimo zaznamenává Codecademy streak uživatele, neboli počet dní, po které se uživatel alespoň jednou denně učil bez přerušení.

% Final
\paragraph{Sledování pokroku dovedností} Unikátní funkcí Codecademy je sledování pokroku u dané dovednosti. Když se uživatel například učí jazyk JavaScript, zvyšuje se mu dovednost vývoje webových aplikací. Uživatel si pak může u každé dovednosti zobrazit jeho úroveň, jak dobře danou látku ovládá a co by se měl například ještě naučit.

% Final
Celkově tak Codecademy obsahuje několik způsobů motivace uživatelů, nicméně zdá se, že má v aplikaci gamifikace spíše vedlejší roli.

% Final
\subsubsection{Placená verze}

% Final
Codecademy, na rozdíl od ostatních aplikací, neobsahuje žádné reklamy. Místo toho nabízí uživatelům několik různých modelů předplatného. Předplatné se dá rozdělit do kategorie pro samouky, studenty a firmy.

% Individual tier
% Final
V rámci kategorie předplatného pro samouky Codecademy nabízí uživatelům některé základní kurzy zdarma. V rámci předplatného Plus se pak mohou uživatelé učit neomezené množství kurzů a dostanou přístup k cestám schopností. Dále mohou neomezeně využívat AI asistenta při učení. Poslední předplatné Pro umožňuje uživatelům se učit pomocí kariérní cesty, získávat profesní certifikáty. Uživatelé s tímto předplatným mohou využívat další funkce aplikace, jako například simulátor pracovních pohovorů, programátorské výzvy a další.

% Premium tier
% Final
Codecademy dále nabízí předplatné pro studenty. Toto předplatné je takřka stejné jako Pro předplatné pro samouky. Hlavním rozdílem je snížená cena, pokud uživatel doloží jeho status studenta.  

% Free tier
% Final
Poslední kategorií předplatného je předplatné pro firmy. Toto předplatné má dvě varianty Teams a Enterprise. Teams nabízí podobné funkce jako Pro předplatné pro samouky s rozdílem, že firma dostane navíc možnost spravovat jednotlivé účty uživatelů, přiřazovat je do skupin a získávat přehledy o učení uživatelů. Enterprise navíc nabízí možnost integrací dalších aplikací a hodiny vedené online instruktorem.

% Final
\subsubsection{Silné a slabé stránky}

% Silné stránky
% Final
Jednou ze silných stránek Codecademy je velké množství výukových materiálů, které stránka nabízí. Uživatelé mají přístup ke kvízům, lekcím, videím, projektům, shrnutím látky a dalším nástrojům, které jim umožní se naučit danou látku.

% Final
Další silnou stránkou Codecademy je minimalistické uživatelské rozhraní. Přestože se některé části uživatelského rozhraní mohou zdát zbytečně složité, aplikace je celkově velice přehledná a snadná na používání, obzvláště pokud se vezme v úvahu množství materiálů, kurzů a funkcí, které aplikace nabízí.

% Slabé stránky
% Final
Jednou z potenciálních nevýhod Codecademy je právě samotné zobrazování vzdělávacích materiálů při učení. Na stránce cvičení má uživatel například přístup k zadání úkolu, vysvětlení látky ve formě textu, dále k videu vysvětlující látku, AI asistentovi, shrnutí látky, komunitní sekci a dalším podpůrným materiálům. Když uživatel prochází jednotlivými cvičeními, může se pak cítit přehlcen tím, že mu jsou všechny informace zobrazovány na jedné stránce.

% Final
Další potenciální nevýhodou Codecademy je nedostatek motivace uživatelů. Aplikace sice obsahuje několik forem gamifikace a motivace uživatele, nicméně zdá se, že tyto způsoby motivace mají v aplikaci spíše vedlejší roli.

% Final
Poslední nevýhodou této aplikace může být to, že je uživatel sám zodpovědný za opakování si dané látky. Pokud například uživatel dokončí kurz, aplikace předpokládá že danou látku uživatel zná a nepomáhá mu si danou látku opakovat. Uživatel tak může například po několika měsících látku zapomenout a pokud by ji následně potřeboval umět, musel by projít kurzem znovu.

% Final
\subsubsection{Další poznatky}

% Final
Kromě výše zmíněných výhod a nevýhod jsem v rámci analýzy narazil na několik funkcí, které stojí za zmínku.

% Zajímavé funkce
% Final
\paragraph{Formátování kódu v editoru} V rámci cvičení má uživatel k dispozici tlačítko, kterým může automaticky zformátovat napsaný kód v editoru. To může být velice užitečné a přispívat ke kvalitní uživatelské zkušenosti, protože uživatel nemusí kód formátovat manuálně.

% Final
\paragraph{Studijní plán} Codecademy umožňuje uživateli si vygenerovat studijní plán. V rámci plánu si může nastavit, jak dlouho a kolik dní v týdnu se chce učit. Aplikace mu pak každý den doporučí konkrétní lekce a aktivity, které má daný den projít. Limitací této funkce je, že si uživatel může sestavit plán pouze pro daný kurz a nemá možnost si vygenerovat plán, který by obsahoval více kurzů zároveň. Přesto může být tato funkce pro uživatele velice užitečná a může značně zjednodušit každodenní učení a mít tak pozitivní vliv na uživatelskou zkušenost.
